% !PRINTED none

\begin{center}
\gesperrt{Thirteenth Chapter.}

\medskip
\gesperrt{General Theory of the Equation of the Fifth Degree.}

\medskip
\rule{0.3\linewidth}{0.4pt}

\medskip
\S.\,106.

\gesperrt{Statement of the Question.}
\end{center}

We have seen in \S.\,44 that beyond the question of the
\gesperrt{Galois} group of an algebraic equation there is a further-reaching
problem, namely the question of the linear substitution group
of the smallest possible number of dimensions to whose form
problem the given equation can be reduced. For the metacyclic
equations, and so in particular for the general equations of the
3rd and 4th degree, this number of dimensions is equal to 1,
which is just what is expressed by saying that these equations
are solvable by radicals. The equations next to be investigated
are then those of the 5th degree, and it will appear that the
solution of the general equation of the 5th degree leads to a
binary linear substitution group, namely to the icosahedron
problem.

But connected with this there is yet another question.
The general equation of the 5th degree contains five coefficients,
which can be regarded as independent variables, and
consequently the root of such an equation is an algebraic
function of five variables. Now one can, however, already
diminish the number of these variables by the simple linear
substitutions. By the \gesperrt{Tschirnhausen} transformation, e.\,g.
to the \gesperrt{Jerrard} form, one can even make the equation depend
on only one variable coefficient (one \gesperrt{parameter}), and one can
thus, through the mediation of equations which do not attain
the 5th degree, make the root of a general equation of the
5th degree depend on an algebraic function of one variable.
